% reconfiguration.tex -- Methodology 3.4: reconfiguration (rewritten 12 Oct 2026 to the thesis plan of 11 Oct:
% notice, re-plan, redistribute, fall back; the three-equation statics; the support margin with its qualifier; the
% constants table; the admission contract; the state machine). 17 Oct (v3): the join is noticed by the co-motion
% confirmation (the default since 14 Oct, proof_rules comotion; card docs/experiments/2026-10-15_attach_comotion_rig.md
% rev 4, Tejen's review); the old contact-and-lift proof is history. Plates clockwise (angle = 30 x plate).
% Code: the planner's detach mode (detection, resize), drone_magnet (contact proof, attachment manager, radio mux),
% the take-back in the tracker; the collaborator's stack carries the lift from contact.

\subsection{Reconfiguration}
\label{sec:reconfig}

A reconfiguration changes the number of UAVs carrying the ring in the air. Both directions are handled by the same
four steps: the planner \emph{notices} the change from an observation, \emph{re-plans} by switching to a solver
built for the new team, \emph{redistributes} the ring's weight across it, and \emph{falls back} to a simple exit
when a step fails. Figure~\ref{fig:reconfig_states} shows the two lanes and their exits; the constants every step
uses are in Table~\ref{tab:constants} (Appendix~\ref{sec:app_reconfig}), each with the reason for its value.

\begin{figure}[htbp]
\centering
\begin{tikzpicture}[>=Stealth, font=\small,
    st/.style={draw, thick, rounded corners=3pt, align=center, minimum height=0.8cm, minimum width=2.2cm, font=\small},
    ours/.style={st, fill=blue!6, draw=blue!50!black},
    his/.style={st, fill=orange!10, draw=orange!60!black},
    exit/.style={st, fill=gray!14, draw=gray!60},
    ev/.style={font=\scriptsize, align=center}]
  % leave lane (top)
  \node[ours] (carry)   at (0,0)      {Carry, $N$ UAVs\\ \scriptsize planner owns the ring};
  \node[ours] (resize)  at (4.8,0)    {Resize to $N-1$\\ \scriptsize warm start, free tensions};
  \node[ours] (carry3)  at (9.6,0)   {Carry, $N-1$\\ \scriptsize statics sets the shares};
  \node[ours] (freed)   at (4.8,-1.7) {Freed UAV\\ \scriptsize steps out 0.5\,m};
  \node[exit] (dland)   at (9.6,-1.7){Lands by itself\\ \scriptsize disarmed alone at rest};
  \draw[->, thick] (carry) -- node[ev, above] {DETACH, or\\ rod stretch} (resize);
  \draw[->, thick] (resize) -- (carry3);
  \draw[->, thick] (resize) -- node[ev, right] {the leaver} (freed);
  \draw[->, thick] (freed) -- (dland);
  % join lane (middle)
  \node[his]  (mission) at (0,-3.6)   {Newcomer, free\\ \scriptsize collaborator's controller};
  \node[his]  (contact) at (4.8,-3.6) {Descent to contact\\ \scriptsize then a small lift};
  \node[ours] (proof)   at (9.6,-3.6){Confirmed by mocap\\ \scriptsize tip moves with the plate};
  \node[ours] (takeback) at (4.8,-5.3){Take-back\\ \scriptsize moved clear, lands alone};
  \node[ours] (fold)    at (9.6,-5.3){Mux switch, fold in\\ \scriptsize seed 0.1\,N, resize to $N+1$};
  \draw[->, thick] (mission) -- node[ev, above] {transit\\ to the plate} (contact);
  \draw[->, thick] (contact) -- node[ev, above] {contact} (proof);
  \draw[->, thick] (proof) -- node[ev, right] {confirmed} (fold);
  \draw[->, thick, dashed] (contact) -- node[ev, right] {give-up} (takeback);
  \draw[->, thick, dashed] (proof.south west) -- node[ev, sloped, above, pos=0.45] {not confirmed} (takeback.north east);
  \draw[->, thick] (fold.east) -- ++(0.55,0) |- node[ev, pos=0.25, right] {carry, $N+1$} (carry3.east);
  % exits (bottom)
  \node[exit] (drop) at (0,-7.0)    {Drop and land\\ \scriptsize second loss, or tilt $>$ 30$^\circ$};
  \node[exit] (cut)  at (4.8,-7.0)  {Tilt cut\\ \scriptsize newcomer's motors off};
  \node[exit] (land) at (9.6,-7.0) {LAND\\ \scriptsize every UAV, from any phase};
  \draw[->, thick, dashed] (takeback) -- node[ev, right] {stuck, tilting} (cut);
\end{tikzpicture}
\caption[Leave and join lanes and their exits]{The leave lane (top), the join lane (middle) and the exits
(bottom). Blue: this work's planner and trackers; orange: the collaborator's controller; grey: an exit that ends
with every UAV owned and on the floor. A leave is noticed from the DETACH command or from the leaver's rod reading
longer than itself; a join from a contact and, after the collaborator's lift, the magnet tip seen moving with the
plate on a taut tether, from motion capture alone. Every exit is reachable from every state: LAND and the drop rule
are not drawn into each.}
\label{fig:reconfig_states}
\end{figure}

\subsubsection{Notice}

\paragraph{A leave.} An \emph{announced} leave is the DETACH command: the leaver's magnet is switched off on its
radio channel and the planner resizes on the same tick, or, for a \emph{ramped} release, its planned tension is run
down over a short ramp before the magnet is released, so that the survivors take up the share before the rod goes.
An \emph{unannounced} leave is a magnet switched off without the planner being told. The planner notices it from the
rod: a rigid rod on its plate cannot measure longer than itself, so when the leaver's rod reads more than 0.06\,m over
its length for two consecutive ticks, the planner drops that UAV from its problem, with no force sensor and no message
from the leaver.

\paragraph{A join.} A join is noticed in two steps, both read from motion capture and nothing else, with the
magnet tip expressed in the plate's own frame. A \emph{contact} is the tip inside the catch radius of the plate
point (24\,mm on the rig), between 15\,mm below and 12\,mm above the plate's face, and slower than 0.03\,m/s for
0.3\,s. A contact is only a candidate: a tip resting on the plate and a magnet that has caught look alike. The
\emph{confirmation} uses the collaborator's small lift: 1.0\,s after its controller signals the lift done, the catch
is confirmed if the tip has stayed within 6\,mm of its contact point while the plate has moved with it (a range of
5\,mm or a rise of 3\,mm), and the tether reads taut against the free hang measured on the same flight (within
3\,mm). A separation of 45\,mm for 0.12\,s fails it, and no hold spans missing data: a gap of more than 0.1\,s
restarts the contact and the confirmation. In the twin, whose stiff ring hardly moves and whose magnet catches at the
instant of contact, the plate-motion condition is waived. This co-motion confirmation replaced, on 14 October, a proof
that rejected the rig's one catch on its taut check.

\subsubsection{Re-plan}

The attachment points $\boldsymbol{\rho}_i$ and rod lengths $l_i$ are runtime parameters of the model, so only the team
size fixes the problem's dimensions, and one solver per team size, built in advance and primed on the measured state,
covers any layout. At the event the planner swaps to the solver for the new size, keeps the survivors on their actual
plates (for a newcomer, the plate it caught), and warm-starts the new problem from the measured state with the tensions
and rod directions free at its first node: the first solve redistributes the load, with no separate hand-off phase. An
announced leave is refused during LAND and when it would leave fewer than three UAVs, a pendulum on two rim points; an
unannounced loss that leaves fewer than three is a fault, which drops the ring (Section~\ref{sec:safety}).

\subsubsection{Redistribute}

\paragraph{What statics fixes.} With three carriers the vertical shares are not the planner's to choose. Take the
ring's centre of mass as the origin of its frame and the survivors' plates at $\boldsymbol{\rho}_i = r(\cos\alpha_i,
\sin\alpha_i, 0)$; with $f_i$ the vertical component of rod $i$'s pull on a level ring, the force balance and the
two moment balances about the horizontal axes are three equations in three unknowns,
\begin{equation}
    \sum_{i=1}^{3} f_i = m_L g, \qquad \sum_{i=1}^{3} f_i \cos\alpha_i = 0, \qquad \sum_{i=1}^{3} f_i \sin\alpha_i = 0,
    \label{eq:statics}
\end{equation}
whose solution is unique whenever the three plates do not lie on a line: each share is the weight times the
barycentric coordinate of the centre of mass in the survivors' triangle. The rods' horizontal components, which the
planner does set through the rod directions, balance one another and make no moment about the centre of mass when
the centre lies in the plates' plane \cite{michael2011cooperative, tognon2018internal}; on the
ring the plates sit 2\,cm above the centre of mass, so the shares of \eqref{eq:statics} hold to within a few per
cent. After a leave the planner's first solve therefore hands the leaver's share to the survivors as
\eqref{eq:statics} says (Figure~\ref{fig:detach}, Appendix~\ref{sec:app_secondary}). A level ring is possible only when every $f_i$ is positive, which
is the condition of the next paragraph.

\paragraph{Support margin.} A rigid payload can hang level from three or more rim points only if its centre of mass
lies inside the polygon of the attachment points \cite{michael2011cooperative}. For plates on a ring of radius $r$
the nearest edge of the survivors' polygon is the chord across their largest angular gap $g_{\max}$, and the
distance from the centre to that chord,
\begin{equation}
    m = r \cos\!\big(g_{\max}/2\big),
    \label{eq:margin}
\end{equation}
is the support margin: positive inside, zero on the edge, negative outside. On the 0.225\,m plate radius a gap of
120, 135, 150, 165 or 180$^\circ$ leaves a margin of 112, 86, 58, 29 or 0\,mm. Removing one carrier from an even
ring of $N$ leaves a largest gap of $4\pi/N$, which keeps the centre inside only for $N \ge 5$; an even ring of
four sits exactly on the boundary after any loss, which is why the four-UAV flights use plates 3, 7, 9 and 11,
where the loss of the carrier on plate 9 leaves an even triangle with a 112\,mm margin. The margin is the lever the
results chapter varies in the twin (Section~\ref{sec:res_leave}).

\paragraph{Which carrier may leave, and what control still sets.} The condition also says which carriers a layout
can afford to lose, and it can be checked before a flight from the plate angles alone. Removing one carrier merges
the two angular gaps on either side of it, and the four gaps of a four-carrier layout sum to $360^\circ$, so some
pair of adjacent gaps sums to at least $180^\circ$: no four-carrier layout on the ring can lose an arbitrary
carrier and still hang level, and a layout admits at most two leavers. The carry layout admits exactly one, the
carrier on plate 9; losing the carrier on plate 7 or plate 11 puts the centre on an edge, and losing the one on plate
3 puts it 112\,mm outside. The layout on plates 3, 6, 9 and 11 admits the same leaver with a $150^\circ$ gap and a
58\,mm margin. A team that must tolerate the loss of any one carrier therefore needs five, and a four-carrier layout
is designed for a known leaver, which is what the rejoin mission does.
% The numbers: tools/admissibility_map.py (results/analysis/admissibility/), from metrics.cog_margin and the
% planner's balanced_tensions; 0 of the 495 four-plate layouts admits every leaver (11 Oct).
With three survivors no controller can change how the
ring's weight is shared; what the planner still sets is the margin, through the layout and the leaver chosen before
the flight, when and how fast the survivors take up the step, and the rods' directions.

\paragraph{Fall back.} After a release the freed UAV steps 0.5\,m out and lands by itself, and is disarmed alone at
its rest; a newcomer whose catch is not confirmed is taken back by this work's tracker and landed; the drop rule, the
tilt cut and LAND are the exits of Section~\ref{sec:safety}. An announced release can be ramped, and a newcomer is
seeded at 0.1\,N and blended to its share. Appendix~\ref{sec:app_reconfig} gives these steps in detail and the constants
of every step with their reasons (Table~\ref{tab:constants}).

\subsubsection{The admission contract}
\label{sec:admission}

The join is a contract between two controllers, written so that nothing on the newcomer's side has to be trusted. The
collaborator's controller must descend onto its plate to contact with the magnet on, lift by a small set amount and
hold, signal that the lift is done, and give up its authority when the catch is confirmed, or give up and go home. This
work must observe the contact and the lift from motion capture only, switch the newcomer's radio to its own tracker on
the tick of the confirmation (through a multiplexer that switches only onto an armed command above idle throttle, and
latches), fold the newcomer in from 0.1\,N, and take it back if nothing is confirmed. The full mission, in which the
released UAV flies the collaborator's errand and rejoins, is built on the same contract (Appendix~\ref{sec:app_reconfig}).
